\documentclass[12pt]{article}
% TikZ setup
\usepackage{tikz}
\usetikzlibrary{arrows.meta,positioning,bending}
\newcommand{\duedate}{11/21/2025}
\newcommand{\assignment}{Assignment 2} % Change to "Problem Set X"

% Change the following to your name and UNI.
\newcommand{\name}{Aksel Kretsinger-Walters, adk2164}
\newcommand{\email}{adk2164@columbia.edu}

% NOTE: Defining collaborators is optional; to not list collaborators, comment out the
% line below. Maximum of two collaborators per problem set.
%\newcommand{\collaborators}{Vig Vigerton (\texttt{UNI}), Alice Bob (\texttt{UNI})}
% No collaborators on PS 0

\makeatletter
\def\input@path{{../}{../../}} % add as many parents as you need
\makeatother
\input{pset_template.tex} %% DO NOT CHANGE THIS LINE

\begin{document}

\psetheader %% DO NOT CHANGE THIS LINE

\section*{(1) Written Problems: Batch Size vs. Learning Rate}
\paragraph{(a)}As batch size increases, how do you expect the optimal learning rate to change?

As batch size increases, I expect the optimal learning rate to also increase. This is because larger batch
sizes provide a more accurate estimate of the gradient, allowing for larger steps to be taken in the

\subsection*{1.1 Training steps vs. batch size}

\paragraph{(a)} For the three points $(A,B,B)$, in Figure 1, which one has the most efficient batch size.
Point C has the most efficient batch size. Lets assume that point C is the batch size s.t. the compute is
used quasi-optimally. This means that as batch size drastically decreases, the number of batches needed to
cover the entire dataset surpasses the number of available compute units, leading to underutilization of
resources (point A resides in this region of the chart - many more training steps are required). On the other
hand, point B's batch sizes are so large that they exceed the capacity of the compute units, leading to
inefficiencies due to memory constraints and increased communication overhead.

\paragraph{(b)} Fill in the blanks.

\tbf{Point A: } Regime: Noise-dominated \quad Potential acceleration: Increase batch size
\tbf{Point B: } Regime: Curvature-dominated \quad Potential acceleration: Decrease batch size

\subsection*{(1.2) Model size, dataset size, and compute}
\paragraph{(a)} Best way to improve model \tbf{test} performance
(C) Increase the model size. With more PF-days of compute, increasing the model size allows for capturing
more complex patterns in the data, leading to better generalization and improved test performance.
\end{document}
